% Cortes mecánicos íntegros: parte_iv.tex:55--106,337--408.
% Residencia material del ensamblaje sucesor de 102 capítulos.
% Residencia de realización: capítulo 76.
\section{Suspensión espacial}

Sea
\[
 \Gamma(\theta)=
 \bigl(r(\theta)\cos\theta,r(\theta)\sin\theta,z(\theta)\bigr).
\]
Su curvatura y torsión se calculan mediante
\[
 \kappa_{\mathrm{F}}=\frac{\|\Gamma'\times\Gamma''\|}{\|\Gamma'\|^3},
 \qquad
 \mathcal T_{\mathrm{F}}=
 \frac{\det(\Gamma',\Gamma'',\Gamma''')}
 {\|\Gamma'\times\Gamma''\|^2}.
\]
Una hélice general en el sentido de Lancret satisface
\[
 \frac{\mathcal T_{\mathrm{F}}}{\kappa_{\mathrm{F}}}=\text{constante}.
\]
Una suspensión radial arbitraria no tiene por qué cumplir esa condición. El
nombre \emph{suspensión helicoidal radial} conserva la estructura de
enrollamiento sin afirmar una propiedad diferencial no demostrada.

\section{Tricromía geométrica}

Los generadores
\[
 J_{+1}^2=-I,\qquad J_0^2=0,\qquad J_{-1}^2=I
\]
producen evoluciones elíptica, parabólica e hiperbólica. Una historia puede
alternarlas y publicar una curva cuya curvatura de Frenet varía de manera
continua. La \emph{tricromía} describe la secuencia de regímenes internos;
la geometría de Frenet describe la salida. El lector radial conserva ambas
sin confundirlas.

\begin{figure}[htbp]
\centering
\begin{tikzpicture}[>=Latex,node distance=10mm and 16mm,
 n/.style={draw=HMTNavy,fill=white,text width=34mm,minimum height=10mm,align=center,font=\small\sffamily}]
\node[n] (int) {historia TRIT\\\(+1,0,-1,\ldots\)};
\node[n,right=of int] (rad) {lector radial\\\(p,\Lambda,C,\mu\)};
\node[n,right=of rad] (curve) {curva publicada\\\(r(\theta),z(\theta)\)};
\node[n,below=of curve] (frenet) {geometría de Frenet\\\(\kappa_{\mathrm{F}},\mathcal T_{\mathrm{F}}\)};
\draw[->,HMTBlue,thick] (int)--(rad);
\draw[->,HMTBlue,thick] (rad)--(curve);
\draw[->,HMTBlue,thick] (curve)--(frenet);
\draw[->,HMTRed,dashed,bend left=28] (frenet.west) to node[below,font=\scriptsize\sffamily]{no selecciona retrospectivamente} (int.south);
\end{tikzpicture}
\caption{Separación causal entre régimen interno y curvatura visible.}
\label{espiral:fig:trit-frenet}
\end{figure}

\chapter{Estructura dinámica del espacio y consecuencias}

\section{Laminación de hojas helicoidales}

La síntesis geométrica es una laminación cuyos elementos son historias
levantadas. Cada hoja porta:
\begin{enumerate}
\item una fase angular;
\item una coordenada radial interna;
\item una profundidad de memoria;
\item una secuencia de hojas APP;
\item un régimen TRIT local;
\item una holonomía TPK y su frontera.
\end{enumerate}
Las secciones transversales son círculos o curvas de nivel; las familias de
secciones forman cilindros; la variación radial produce suspensiones; la
involución produce hojas conjugadas. La geometría real es la imagen de esta
laminación bajo una publicación que puede contraer varias hojas.

\section{Clasificación cinemática por incrementos angular y radial}

Una recta radial aparece cuando \(\omega=0\) y \(\beta\neq0\). Una
circunferencia aparece cuando \(\omega\neq0\) y \(\beta=0\). La espiral
aparece cuando ambos son no nulos. Estas tres salidas son configuraciones de
un mismo mecanismo de transporte, no objetos sin relación. El soporte APP
es bidimensional desde el comienzo; el recorrido unidimensional es una
sección particular de ese soporte.

\section{Cociente de Möbius mediante retorno con inversión de orientación}

Cuando la involución de ruta se identifica en la costura después de una
vuelta, el cociente puede adquirir una banda de Möbius. La afirmación exige
una acción libre y el entrelazador que transporte la orientación. La banda no
se deduce de la mera presencia de un signo \(-1\); se deduce de la costura
\[
 (x,0)\sim(C_{\mathrm{rev}}x,1)
\]
en el dominio declarado. En el atlas material, la involución dodecafásica y
la inversión celular continúan siendo operaciones distintas.

\Needspace{34\baselineskip}
\section{Clasificación de 3+2 puntos de equilibrio mediante el Hessiano}

La intuición de los tres puntos colineales y los dos triangulares de Lagrange
puede formularse como una prueba posterior. Dado un potencial efectivo
\(V_{\mathrm{eff}}\) producido por la realización HMT--MD, los puntos críticos
satisfacen
\[
 \nabla V_{\mathrm{eff}}=0.
\]
Su clasificación exige el Hessiano y el índice de Morse. Sólo si el mapa
dimensional produce exactamente tres críticos colineales y dos triangulares
puede relacionarse con la partición \(3+2\). La correspondencia no se usa
como entrada de la taxonomía espiral.

\section{Jerarquía causal entre curvas, campos y secciones materiales}

La correspondencia espiral abre tres salidas coordinadas:
\[
 \text{ledger APP--TRIT--TPK}
 \longrightarrow
 \begin{cases}
 \text{curva arquimediana},\\
 \text{modo electromagnético},\\
 \text{firma de ruta material}.
 \end{cases}
\]
La primera está cerrada algebraicamente en este trabajo. La segunda posee
una realización exacta para modos circulares y una condición de
compatibilidad espectral. La tercera dispone del electrón central y del
atlas, pero requiere verificar el cuadrado de constancia sobre cada
multisección antes de asignar taxonomías de sabor.
